\paragraph{Theorem F.2.27 (Proved: the entire remaining order-six locus is excluded).}
Every simply connected datum in the second row of (F.S1),
at every rational relative position and every stable integral
lattice with primitive maximal cusp, has
\[
                 \operatorname{Tor}_3 H_3(M;\mathbb Z)\ne0
\]
if its degree-two suspension cokernels vanish.
Consequently no such datum is three-connected.  Together with
F.2.25 this leaves only the minimal order pair $(5,4)$ as a
possible three-connected datum in the full problem F.2.
\begin{proof}
We use cohomology with coefficients in the discrete valuation
ring $R=\mathbb Z_3$.  Suppose the degree-two suspension
cokernels vanish; otherwise F.2.25 already gives nonzero
$H_2$ or $H_3$.  That theorem forces $m=n=h=6$ at the
order-six point, in the notation (F.S3)--(F.S5).
We first determine its local integral representation over $R$.

The splitting by the fixed vector $u$ is integral:
$\Lambda_R=Ru\oplus H_R$.  On $H_R$, the finite map $D$
has characteristic polynomial
$(x-1)(x+1)^2\Phi_3(x)$.  Since two is a unit in $R$,
the idempotents $(I\pm D^3)/2$ give an integral direct sum
\[
 H_R=L_+\oplus L_-,\qquad
 \operatorname{rank}L_+=3,\quad
 \operatorname{rank}L_-=2.
\]
On $L_-$ the map is exactly $-I$, by its rational spectrum.
On $L_+$ it has order three.  The primitive fixed vector
$\lambda$ lies in $L_+$, and the primitive invariant
covector $\varphi$ vanishes on $L_-$, so their restrictions
remain primitive.  Equation (F.S5) with $h=6$ gives
\[
 2(I+D+D^2)|_{L_+}=\lambda\varphi|_{L_+}.
 \tag{F.X1}
\]
This operator is nonzero modulo three.  But modulo three
$I+D+D^2=(D-I)^2$ and $(D-I)^3=0$.  The reduction of $D$
on the rank-three space therefore has one Jordan block
of length three.  Choose $q\in L_+$ whose image has
$(D-I)^2q\ne0$ modulo three.  Then $q,Dq,D^2q$ are a
basis modulo three and hence an $R$-basis.  In this basis
$D$ is the three-cycle permutation matrix, since $D^3=I$.
We have proved, without a global lattice classification,
that the full finite representation over $R$ is
\[
                   1\oplus\operatorname{Perm}_3\oplus(-I_2).
 \tag{F.X2}
\]
Its dual has the same form, by reversing the three-cycle.

Let $u,a,b,c,d,e$ be a cohomology basis for (F.X2), with
$a,b,c$ cycled and $d,e$ negated.  A full integral basis
of the invariant degree-three lattice is
\[
 u(ab+bc+ca),\qquad abc,\qquad ude,\qquad(a+b+c)de.
 \tag{F.X3}
\]
Here juxtaposition denotes exterior multiplication.
To check completeness, terms with an odd number of
$d,e$ are negated by the cube of the finite map and
have no invariants.  The terms with zero or two such
factors give exactly the four orbit sums in (F.X3),
each primitive.  With volume $uabcde$, their skew
pairing has one block with coefficient three and one
block with coefficient a unit.  Thus its determinant
has three-adic valuation two.  The exact certificate
\texttt{fourfolds\_universal\_56\_obstruction.py} checks
the full invariant kernel, its primitivity, and the
pairing independently by exterior multiplication and
six-column determinants; its Smith factors are
$(1,1,3,3)$.

Let $B_6$ be the actual reduced finite fibre.  Its torus
cover has degree six.  Its free middle cohomology has
rank four and unimodular skew pairing, by integral
Poincare duality.  If $I$ denotes the index of its
pullback lattice in the invariant lattice, multiplication
of pairings by the covering degree gives
\[
 2v_3(I)+2=4v_3(6)=4,
 \qquad v_3(I)=1.
 \tag{F.X4}
\]
Consequently the three-primary degree-three covering
cokernel is $R/3R$.  Translations of the affine action
do not enter this calculation: the cover degree,
linear invariant lattice and oriented pairings are
those of the actual free quotient for every allowed
affine numerator.

We explain why this local group must appear in total
homology even when the global integral lattice varies.
Let $\mathcal V$ be the rational rank-six cohomology
local system.  It is an extension of the irreducible
rank-five system $W$ by a trivial line.  The determinant
of $W$ is trivial.  Irreducibility in odd dimension
implies that neither $\bigwedge^2 W$ nor its dual
has a nonzero invariant: an invariant alternating
form would have an invariant radical and would have
to be nondegenerate, impossible in odd dimension.
Also $\bigwedge^3W\simeq\bigwedge^2W^*$ by its
trivial determinant.  The exterior filtration therefore
gives
\[
 (\bigwedge^3\mathcal V)^\Gamma=0,\qquad
 ((\bigwedge^2\mathcal V)^*)^\Gamma=0.
 \tag{F.X5}
\]
These rational vanishing statements remain true over
$\mathbb Q_3$, since they are kernels of rational
matrices.

As in (F.S7), take the image of
$R^3g_*R\to j_*\bigwedge^3\mathcal V_R$.
Its kernel and cokernel are point-supported sheaves.
The first vanishing in (F.X5) makes the local group
$R/3R$ inject into $H^1(B,R^3g_*R)$.  Bidegree $(1,3)$
has no incoming or outgoing Leray differential on
the two-dimensional base, so this injection survives
to $E_\infty^{1,3}$.

There is one extension issue, and it cannot remove
this torsion.  The second vanishing in (F.X5) says
that $H^2(B,R^2g_*R)$ has rational rank zero.
Point-supported changes do not alter this degree-two
base cohomology; it is the coinvariant module, as
in (E.9).  It is a finitely generated torsion $R$-module
and hence a finite group.  Therefore
$E_\infty^{2,2}$ is finite as well.  In the filtration
of $H^4(M;R)$, take the inverse image of the injected
$R/3R\subset E_\infty^{1,3}$.  It is an extension
of $R/3R$ by the finite group $E_\infty^{2,2}$.
It is thus a nonzero finite three-primary subgroup
of $H^4(M;R)$; in particular it cannot be absorbed
as a finite quotient of a free filtration kernel.
Universal coefficients identify this torsion with
the three-primary torsion of $H_3(M;\mathbb Z)$.
This proves the claim at every rational position and
every integral lattice, and completes the exclusion.
\end{proof}
